% Generated draft: requires mathematical and source review.
\documentclass[11pt]{article}
\usepackage[margin=25mm]{geometry}
\usepackage{fontspec}
\setmainfont{Latin Modern Roman}
\usepackage{amsmath,amssymb,mathtools,unicode-math}
\setmathfont{Latin Modern Math}
\usepackage{enumitem,needspace,xurl,hyperref}
\setlength{\parindent}{0pt}
\setlength{\parskip}{5pt}
\setlength{\emergencystretch}{3em}
\newcommand{\N}{\mathbb{N}}
\newcommand{\Z}{\mathbb{Z}}
\newcommand{\Q}{\mathbb{Q}}
\newcommand{\R}{\mathbb{R}}
\newcommand{\C}{\mathbb{C}}
\newcommand{\F}{\mathbb{F}}
\newcommand{\A}{\mathbb{A}}
\newcommand{\PP}{\mathbb{P}}
\newcommand{\E}{\mathbb{E}}
\newcommand{\eps}{\varepsilon}
\newcommand{\dd}{\,\mathrm{d}}
\begin{document}
\section*{UnsolvedMath ID: 700403}
\textbf{Source problem code:} catalogue entry
\subsection{Definitions and mathematical statement}
Let \(G\) be an abstract group with identity \(1\). The group is \emph{periodic} if every element has finite order, that is,
\[
\forall g\in G\;\exists n\in\Z_{\geq 1}\quad g^n=1.
\]
An \emph{involution} is an element of order exactly \(2\); thus \(G\) has no involutions if
\[
\forall g\in G\quad (g^2=1\Longrightarrow g=1).
\]
An infinite group is \emph{non-simple} precisely when it has a normal subgroup \(N\) with \(\{1\}\subsetneq N\subsetneq G\).

The supplied record does not define ``biprimitively finite'' or the parenthetical variant ``(conjugacy) biprimitively finite.'' Write \(\mathcal P_{\mathrm b}(G)\) and \(\mathcal P_{\mathrm c}(G)\) as placeholders for these respective intended properties. Their defining criteria require confirmation from the cited source; they are not supplied mathematical definitions.

For each intended choice \(\mathcal P\) among these properties, the corresponding question is whether, for every group \(G\),
\[
\begin{gathered}
G\text{ is infinite and periodic},\qquad
\mathcal P(G),\qquad
\forall g\in G\;(g^2=1\Longrightarrow g=1)
\end{gathered}
\quad\Longrightarrow\quad
\exists N\trianglelefteq G\quad \{1\}\subsetneq N\subsetneq G.
\]
This preserves the structural question, but a complete intrinsic formulation remains dependent on the missing definitions and clarification of the parenthetical qualification.
\subsection{Short English statement}
Must every infinite periodic group without involutions that satisfies the intended ``(conjugacy) biprimitively finite'' condition have a nontrivial proper normal subgroup? The supplied record leaves that condition undefined and its parenthetical scope unresolved.
\subsection{Sources}
\begin{itemize}
\item Local UnsolvedMath registry, permanent ID \(700403\), exact source problem code \texttt{catalogue entry}, title ``Kourovka Notebook Problem 6.61.'' Supplied upstream revision: \texttt{d371ee917ac9647fbb06cd6fc529deec6f3dda82}; record updated \texttt{2026-10-16T14:00:00Z}.
\item Supplied attribution: V. P. Shunkov, proposed in 1978; \emph{The Kourovka Notebook (Unsolved Problems in Group Theory)}, 6th issue (1978), Problem 6.61. The record references the 21st issue edition (2026), \texttt{arXiv:1401.0300v48}. The supplied collection description names E. I. Khukhro and V. D. Mazurov as editors.
\item The registry labels this problem ``open.'' This is supplied source metadata, not an independently verified assessment. The cited edition and its terminology have not been independently checked for this draft.
\end{itemize}
\end{document}
